\documentclass[aspectratio=169,11pt]{beamer}

\usetheme{AnnArbor}
\usecolortheme{dolphin}
\usefonttheme{professionalfonts}
\setbeamertemplate{navigation symbols}{}
\setbeamertemplate{caption}[numbered]

\usepackage{fontspec}
\setsansfont{DejaVu Sans}
\setmonofont{DejaVu Sans Mono}
\usepackage{amsmath,amssymb}
\usepackage{booktabs}
\usepackage{array}
\usepackage{tikz}
\usetikzlibrary{arrows.meta,positioning,shapes.geometric,fit}
\usepackage{hyperref}
\hypersetup{colorlinks=true,linkcolor=blue,urlcolor=blue}

\title[DBMS]{Hệ quản trị cơ sở dữ liệu và MySQL}
\subtitle{Course Roadmap \& Learning Overview}
\author{Vu Duc Minh}
\institute{National Economics University, Hanoi, Vietnam}
\date{\today}

\definecolor{dbblue}{RGB}{20,74,135}
\definecolor{dblight}{RGB}{232,242,255}
\definecolor{dbgreen}{RGB}{28,120,90}
\definecolor{dborange}{RGB}{210,120,30}

\setbeamercolor{structure}{fg=dbblue}
\setbeamercolor{title}{fg=white,bg=dbblue}
\setbeamercolor{frametitle}{fg=dbblue}
\setbeamercolor{block title}{fg=white,bg=dbblue}
\setbeamercolor{block body}{bg=dblight,fg=black}

\newcommand{\kw}[1]{\textcolor{dbblue}{\textbf{#1}}}
\newcommand{\smallgap}{\vspace{0.4em}}

\begin{document}

%------------------------------------------------
\begin{frame}
  \titlepage
\end{frame}

%------------------------------------------------
\begin{frame}{Mục tiêu của học phần}
Sau học phần, người học cần có thể:
\begin{itemize}
  \item Giải thích vai trò của \kw{cơ sở dữ liệu} và \kw{DBMS} trong hệ thống thông tin.
  \item Thiết kế mô hình dữ liệu ở mức khái niệm bằng \kw{ER/EER model}.
  \item Chuyển mô hình ER sang \kw{lược đồ quan hệ} và phân tích khóa, phụ thuộc hàm.
  \item Áp dụng \kw{chuẩn hóa} để giảm dư thừa và bất nhất dữ liệu.
  \item Sử dụng \kw{MySQL} để quản trị, truy vấn, tối ưu và lập trình trong cơ sở dữ liệu.
\end{itemize}
\end{frame}

%------------------------------------------------
\begin{frame}{Bức tranh tổng thể của học phần}
\centering
\begin{tikzpicture}[
  node distance=1.05cm and 1.5cm,
  box/.style={rounded corners, draw=dbblue, thick, fill=dblight, minimum width=3.1cm, minimum height=0.85cm, align=center, font=\small},
  arr/.style={-Latex, thick, draw=dbblue}
]
\node[box] (intro) {Nền tảng DBMS\\kiến trúc, abstraction};
\node[box, right=of intro] (er) {ER / EER Model\\thiết kế khái niệm};
\node[box, right=of er] (rel) {Relational Model\\khóa, FD, schema};
\node[box, below=of rel] (norm) {Normalization\\1NF--4NF};
\node[box, left=of norm] (mysql) {MySQL Server\\cài đặt, vận hành};
\node[box, left=of mysql] (sql) {SQL \& Admin\\query, index, transaction};
\draw[arr] (intro) -- (er);
\draw[arr] (er) -- (rel);
\draw[arr] (rel) -- (norm);
\draw[arr] (norm) -- (mysql);
\draw[arr] (mysql) -- (sql);
\end{tikzpicture}

\smallgap
\textit{Tư duy chính: đi từ mô hình hóa dữ liệu đến triển khai và khai thác dữ liệu bằng MySQL.}
\end{frame}

%------------------------------------------------
\section{Introduction}
\begin{frame}{Phần 1: Introduction to DBMS}
\begin{block}{Nội dung chính}
\begin{itemize}
  \item Cơ sở dữ liệu và hệ quản trị cơ sở dữ liệu.
  \item Nhu cầu sử dụng DBMS trong hệ thống thực tế.
  \item Kiến trúc 1-tier, 2-tier, 3-tier.
  \item Trừu tượng hóa dữ liệu và độc lập dữ liệu.
  \item Schema, instance và mô hình dữ liệu.
\end{itemize}
\end{block}

\begin{alertblock}{Câu hỏi dẫn nhập}
Tại sao một ứng dụng kinh doanh không nên lưu dữ liệu quan trọng bằng các file rời rạc?
\end{alertblock}
\end{frame}

%------------------------------------------------
\begin{frame}{DBMS giải quyết vấn đề gì?}
\begin{columns}[T,onlytextwidth]
\column{0.48\textwidth}
\begin{block}{Khi không dùng DBMS}
\begin{itemize}
  \item Dữ liệu trùng lặp, khó kiểm soát.
  \item Khó đảm bảo tính nhất quán.
  \item Bảo mật và phân quyền thủ công.
  \item Khó sao lưu, phục hồi, mở rộng.
\end{itemize}
\end{block}
\column{0.48\textwidth}
\begin{block}{Khi dùng DBMS}
\begin{itemize}
  \item Quản lý dữ liệu tập trung.
  \item Ràng buộc toàn vẹn dữ liệu.
  \item Truy vấn bằng SQL.
  \item Hỗ trợ giao dịch, bảo mật, backup.
\end{itemize}
\end{block}
\end{columns}
\end{frame}

%------------------------------------------------
\section{ER Model}
\begin{frame}{Phần 2: Entity Relationship Model}
\begin{block}{Vai trò}
ER model giúp mô tả dữ liệu ở mức \kw{khái niệm}, trước khi chuyển thành các bảng trong cơ sở dữ liệu quan hệ.
\end{block}

\begin{columns}[T,onlytextwidth]
\column{0.32\textwidth}
\textbf{Thực thể}\
Khách hàng, đơn hàng, sản phẩm.
\column{0.32\textwidth}
\textbf{Thuộc tính}\
Mã khách hàng, ngày đặt, giá bán.
\column{0.32\textwidth}
\textbf{Mối quan hệ}\
Khách hàng đặt đơn hàng; đơn hàng gồm sản phẩm.
\end{columns}

\smallgap
\begin{alertblock}{Thực hành}
Xác định entities, attributes, keys và relationships cho một hệ thống bán hàng nhỏ.
\end{alertblock}
\end{frame}

%------------------------------------------------
\begin{frame}{Từ bài toán thực tế đến ERD}
\centering
\begin{tikzpicture}[
  node distance=0.9cm,
  process/.style={draw=dbblue, rounded corners, fill=dblight, align=center, minimum width=4.0cm, minimum height=0.82cm, font=\small},
  arr/.style={-Latex, thick, draw=dbblue}
]
\node[process] (req) {Mô tả nghiệp vụ\\yêu cầu dữ liệu};
\node[process, below=of req] (entity) {Xác định thực thể\\và thuộc tính};
\node[process, below=of entity] (rel) {Xác định quan hệ\\và lực lượng tham gia};
\node[process, below=of rel] (erd) {Vẽ ERD\\kiểm tra ràng buộc};
\draw[arr] (req) -- (entity);
\draw[arr] (entity) -- (rel);
\draw[arr] (rel) -- (erd);
\end{tikzpicture}
\end{frame}

%------------------------------------------------
\section{Relational Model}
\begin{frame}{Phần 3: Relational Model and Functional Dependencies}
\begin{block}{Nội dung chính}
\begin{itemize}
  \item Relation, attribute, tuple và relational schema.
  \item Super key, candidate key, primary key, foreign key.
  \item Mapping từ ER model sang relational model.
  \item Functional dependency: determinant và dependent attribute.
  \item Attribute closure và chiến lược thiết kế schema.
\end{itemize}
\end{block}

\begin{exampleblock}{Ví dụ}
Nếu \texttt{StudentID} xác định duy nhất \texttt{StudentName}, ta có phụ thuộc hàm:
\[
\texttt{StudentID} \rightarrow \texttt{StudentName}.
\]
\end{exampleblock}
\end{frame}

%------------------------------------------------
\begin{frame}{Mapping ER sang quan hệ}
\begin{columns}[T,onlytextwidth]
\column{0.52\textwidth}
\begin{itemize}
  \item Mỗi entity mạnh thường trở thành một bảng.
  \item Khóa chính của entity trở thành primary key của bảng.
  \item Quan hệ 1--N thường được biểu diễn bằng foreign key ở phía N.
  \item Quan hệ N--N thường được tách thành bảng trung gian.
\end{itemize}
\column{0.44\textwidth}
\begin{block}{Ví dụ}
\texttt{Customer(CustomerID, Name)}\\
\texttt{Order(OrderID, OrderDate, CustomerID)}
\end{block}
\end{columns}

\begin{alertblock}{Điểm cần kiểm tra}
Một thiết kế tốt phải bảo toàn ý nghĩa nghiệp vụ, khóa và ràng buộc quan hệ.
\end{alertblock}
\end{frame}

%------------------------------------------------
\section{Normalization}
\begin{frame}{Phần 4: Normalization}
\begin{block}{Mục tiêu}
Chuẩn hóa giúp giảm dư thừa dữ liệu và tránh các lỗi cập nhật, chèn, xóa.
\end{block}

\begin{center}
\begin{tabular}{>{\bfseries}p{0.16\textwidth} p{0.72\textwidth}}
\toprule
1NF & Mỗi ô chứa giá trị nguyên tử, không có nhóm lặp.\\
2NF & Không có phụ thuộc bộ phận vào một phần của khóa ghép.\\
3NF & Không có phụ thuộc bắc cầu giữa các thuộc tính không khóa.\\
4NF & Xử lý phụ thuộc đa trị độc lập.\\
\bottomrule
\end{tabular}
\end{center}
\end{frame}

%------------------------------------------------
\begin{frame}{Tư duy chuẩn hóa}
\centering
\begin{tikzpicture}[
  node distance=1.1cm,
  box/.style={draw=dbblue, rounded corners, fill=dblight, align=center, minimum width=3.7cm, minimum height=0.85cm, font=\small},
  arr/.style={-Latex, thick, draw=dbblue}
]
\node[box] (raw) {Bảng ban đầu\\có thể dư thừa};
\node[box, right=of raw] (fd) {Phân tích khóa\\và phụ thuộc hàm};
\node[box, right=of fd] (decomp) {Tách bảng\\theo dạng chuẩn};
\node[box, below=of fd] (check) {Kiểm tra lossless join\\và bảo toàn phụ thuộc};
\draw[arr] (raw) -- (fd);
\draw[arr] (fd) -- (decomp);
\draw[arr] (decomp) -- (check);
\draw[arr] (check) -- (fd);
\end{tikzpicture}
\end{frame}

%------------------------------------------------
\section{MySQL Server}
\begin{frame}{Phần 5: MySQL Server}
\begin{block}{Nội dung triển khai}
\begin{itemize}
  \item Giới thiệu MySQL và hệ sinh thái công cụ.
  \item Cài đặt MySQL Server và MySQL Workbench trên Windows.
  \item Kết nối bằng command line và VS Code.
  \item Start, stop, restart MySQL service.
  \item Nạp cơ sở dữ liệu mẫu \texttt{classicmodels}.
  \item Storage engine, option file, system variables và data directory.
\end{itemize}
\end{block}
\end{frame}

%------------------------------------------------
\begin{frame}[fragile]{Ví dụ thao tác cơ bản với MySQL}
\begin{block}{Tạo và chọn database}
\begin{verbatim}
CREATE DATABASE shop_db
  CHARACTER SET utf8mb4
  COLLATE utf8mb4_unicode_ci;

USE shop_db;
SHOW TABLES;
\end{verbatim}
\end{block}

\begin{alertblock}{Checklist an toàn}
Trước khi \texttt{DROP DATABASE} hoặc \texttt{DROP TABLE}, cần kiểm tra đúng môi trường, đúng tên database và đã có backup nếu dữ liệu quan trọng.
\end{alertblock}
\end{frame}

%------------------------------------------------
\section{SQL}
\begin{frame}{Phần 6: SQL Statements in MySQL}
\begin{columns}[T,onlytextwidth]
\column{0.48\textwidth}
\begin{block}{Data Definition}
\begin{itemize}
  \item Data types.
  \item CREATE TABLE.
  \item Constraints.
  \item ALTER, RENAME, DROP, TRUNCATE.
\end{itemize}
\end{block}
\column{0.48\textwidth}
\begin{block}{Querying Data}
\begin{itemize}
  \item SELECT cơ bản.
  \item JOIN, GROUP BY, HAVING.
  \item Subquery, CTE.
  \item Set operations, window functions.
\end{itemize}
\end{block}
\end{columns}
\end{frame}

%------------------------------------------------
\begin{frame}[fragile]{Ví dụ SELECT trên cơ sở dữ liệu mẫu}
\begin{block}{Tổng doanh thu theo khách hàng}
\begin{verbatim}
SELECT c.customerName,
       SUM(od.quantityOrdered * od.priceEach) AS revenue
FROM customers c
JOIN orders o ON c.customerNumber = o.customerNumber
JOIN orderdetails od ON o.orderNumber = od.orderNumber
GROUP BY c.customerNumber, c.customerName
ORDER BY revenue DESC;
\end{verbatim}
\end{block}
\end{frame}

%------------------------------------------------
\begin{frame}{Index, tối ưu truy vấn và transaction}
\begin{block}{Index Optimization}
\begin{itemize}
  \item Tạo, xóa và xem thông tin index.
  \item Composite index, unique index, prefix index, invisible index.
  \item Đọc \texttt{EXPLAIN} và sử dụng index hints khi cần.
\end{itemize}
\end{block}

\begin{block}{Data Modification \& Transactions}
\begin{itemize}
  \item INSERT, UPDATE, DELETE và thao tác nâng cao.
  \item START TRANSACTION, COMMIT, ROLLBACK, SAVEPOINT.
  \item Locking, row lock, metadata lock và deadlock.
\end{itemize}
\end{block}
\end{frame}

%------------------------------------------------
\section{Programmability}
\begin{frame}{Phần 7: View, Procedure, Function, Trigger, Event}
\begin{block}{Ý tưởng chung}
Các đối tượng này giúp đóng gói logic xử lý trong database, giảm lặp code ở tầng ứng dụng và hỗ trợ tự động hóa một số tác vụ.
\end{block}

\begin{columns}[T,onlytextwidth]
\column{0.5\textwidth}
\begin{itemize}
  \item \kw{View}: đóng gói truy vấn.
  \item \kw{Stored procedure}: đóng gói quy trình xử lý.
  \item \kw{Stored function}: trả về một giá trị.
\end{itemize}
\column{0.5\textwidth}
\begin{itemize}
  \item \kw{Trigger}: tự động chạy khi INSERT/UPDATE/DELETE.
  \item \kw{Event}: tác vụ định kỳ trong MySQL.
\end{itemize}
\end{columns}
\end{frame}

%------------------------------------------------
\begin{frame}{Lộ trình học gợi ý}
\centering
\begin{tikzpicture}[
  node distance=0.65cm,
  week/.style={draw=dbblue, rounded corners, fill=dblight, align=left, minimum width=9.7cm, minimum height=0.62cm, font=\small},
]
\node[week] (w1) {Tuần 1--2: DBMS fundamentals, architecture, abstraction, independence};
\node[week, below=of w1] (w2) {Tuần 3--4: ER/EER modeling và bài thực hành ERD};
\node[week, below=of w2] (w3) {Tuần 5--6: Relational model, keys, ER-to-relational mapping};
\node[week, below=of w3] (w4) {Tuần 7--8: Functional dependencies, closure, normalization};
\node[week, below=of w4] (w5) {Tuần 9--11: MySQL Server, administration, SELECT, JOIN, aggregate};
\node[week, below=of w5] (w6) {Tuần 12--14: Index, transactions, views, procedures, triggers, events};
\end{tikzpicture}
\end{frame}

%------------------------------------------------
\begin{frame}{Hoạt động học tập và đánh giá}
\begin{columns}[T,onlytextwidth]
\column{0.47\textwidth}
\begin{block}{Hoạt động}
\begin{itemize}
  \item Đọc lecture note trước buổi học.
  \item Làm lab trên MySQL và classicmodels.
  \item Vẽ ERD từ tình huống kinh doanh.
  \item Viết truy vấn và giải thích kết quả.
\end{itemize}
\end{block}
\column{0.47\textwidth}
\begin{block}{Sản phẩm đầu ra}
\begin{itemize}
  \item ERD và relational schema.
  \item Script SQL tạo bảng và ràng buộc.
  \item Bộ truy vấn phân tích dữ liệu.
  \item Báo cáo ngắn về thiết kế và tối ưu.
\end{itemize}
\end{block}
\end{columns}
\end{frame}

%------------------------------------------------
\begin{frame}{Tài liệu tham khảo chính}
\begin{itemize}
  \item Carlos Coronel, Steven Morris. \textit{Database Systems: Design, Implementation, and Management}, 14th Edition. Cengage, 2023.
  \item Elvis C. Foster, Shripad Godbole. \textit{Database Systems: A Pragmatic Approach}, 3rd Edition. CRC Press / Taylor \& Francis, 2022.
  \item MySQL Tutorial: tài liệu thực hành về MySQL, SQL và các chủ đề nâng cao.
  \item GeeksforGeeks DBMS and SQL tutorials: tài liệu bổ trợ theo chủ đề.
  \item Roadmap.sh SQL Roadmap: lộ trình tự học SQL.
\end{itemize}
\end{frame}

%------------------------------------------------
\begin{frame}
\centering
\Huge \textbf{Thank you!}

\vspace{1em}
\Large Questions \& Discussion
\end{frame}

\end{document}
